\documentclass[11pt]{article}
\usepackage[a4paper,margin=0.65in]{geometry}
\usepackage[T1]{fontenc}
\usepackage{amsmath,amssymb,amsfonts}
\usepackage{graphicx}
\usepackage{pdfpages}
\usepackage[english,shorthands=off]{babel}
\usepackage{fancyhdr}
\usepackage[bookmarks=false,colorlinks=false]{hyperref}
\usepackage[protrusion=true,expansion=false]{microtype}
\usepackage{array}
\emergencystretch=3em
\tolerance=600
\providecommand{\modd}{\mathrm{mod}}
\providecommand{\Disc}{\mathrm{Disc}}
\providecommand{\canont}{}
\providecommand{\asterism}{*}
\providecommand{\Hessian}{H}
\providecommand{\tome}{}
\providecommand{\page}{p.}
\pagestyle{fancy}\fancyhf{}\fancyhead[L]{Pages 301--325}\fancyhead[R]{Cayley --- Collected Papers, Vol. XIII}\renewcommand{\headrulewidth}{0.4pt}
\setlength{\headheight}{15pt}

\newcommand{\paper}[3]{%
  \clearpage
  \begin{center}
    {\large\bfseries #2}\\[0.5em]
    {\small [#1]}\\[0.3em]
    {\small #3}
  \end{center}
  \vspace{0.5em}
}

\begin{document}

\setcounter{page}{279}

% ====================================================================
% PAPER 932 (continued from prior segment): ON SYMMETRIC FUNCTIONS
% AND SEMINVARIANTS.
% Printed pp.\ 279--303 (PDF pp.\ 301--325).
% (The earlier portion of this paper, printed pp.\ 254--278,
%  corresponds to PDF pp.\ 276--300, which were not produced in
%  any prior segment of Vol.\ XIII.)
% ====================================================================
\begin{center}
{\large\bfseries ON SYMMETRIC FUNCTIONS AND SEMINVARIANTS (continued)}\\[0.4em]
{\small [932, pp.\ 279--303]}
\end{center}
\vspace{0.4em}

% ----- printed page 279 (PDF 301) -----
where the function $U$ operated upon by $\Theta_\sigma$ and $\Theta_{\sigma'}$ respectively is in each case the same function of the coefficients. It is easy to see that, if $T$ is a non-unitary function of the roots, then whatever be the number of the roots we have $S\partial_a \cdot T = a$ determinate symmetric function of the roots, and consequently $=$ a determinate function of the coefficients. We thus have $\Theta_\sigma U$ and $\Theta_{\sigma'} U$ equal to each other; that is,
\[
(\Theta_{\sigma'} - \Theta_\sigma) U = 0;
\]
we may write
\[
\Theta_\sigma = \sigma c \partial_b + (\sigma - 1) b \partial_c + \ldots + \quad p \partial_q,
\]
\[
\Theta_{\sigma'} = \sigma' c \partial_b + (\sigma' - 1) b \partial_c + \ldots + (\sigma' - \sigma + 1) p \partial_q,
\]
for the subsequent terms of $\Theta_{\sigma'}$, as involving $\partial_r$, $\partial_s$, \&c., are inoperative; hence writing
\[
\Delta = \partial_b + b \partial_c + c \partial_d + \ldots + p \partial_q,
\]
or as we may more simply express it
\[
\Delta = \partial_b + b \partial_c + c \partial_d + \ldots,
\]
we have $\Theta_{\sigma'} - \Theta_\sigma = (\sigma' - \sigma) \Delta$, and consequently $\Delta U = 0$; $\Delta$ is thus an annihilator of any function $U$ of the coefficients which is equal to a non-unitary function of the roots; or more shortly $\Delta$ is an annihilator of any non-unitariant.

22.\ Similarly, from the two equations $\Theta_\sigma = -S\partial_a$, and $\Theta_{\sigma'} = S\partial_a$ regarded as operating upon a non-unitary function, we deduce $\sigma' \Theta_\sigma - \sigma \Theta_{\sigma'} = (\sigma - \sigma') S\partial_a$: the left-hand side is here $= (\sigma - \sigma') \Delta_1$, if
\[
\Delta_1 = b \partial_c + 2c \partial_d + 3d \partial_e + \ldots + (\sigma - 1) p \partial_q,
\]
or say
\[
\Delta_1 = b \partial_c + 2c \partial_d + 3d \partial_e + \ldots,
\]
viz.\ we have $\Delta_1 = S\partial_a$; for instance, if as before
\[
T = U = S\alpha^4 = -4e + 4bd + 2c^2 - 4b^2c + b^4,
\]
then
\[
(b\partial_c + 2c\partial_d + 3d\partial_e)(-4e + 4bd + 2c^2 - 4b^2c + b^4) = S\partial_a \cdot S\alpha^4 = 4S\alpha^3 = 4(-3d + 3bc - b^3),
\]
as can be at once verified. It is to be noticed, however, that $S\partial_a$ operating upon a non-unitary function of the roots does not in every case give a non-unitary function; and thus successive operations with $\Delta_1$ will not give a succession of non-unitariants.

23.\ I investigate the foregoing result in regard to $\Delta$ in a different manner; suppose, for instance, that $T = U$ is the non-unitary function $S\alpha^4$ of the roots,
\[
(= -4e + 4bd + 2c^2 - 4b^2c + b^4).
\]
The number of roots is at least $= 4$, and I take it to be $= 4$, say the roots are $\alpha, \beta, \gamma, \delta$. Consider a fifth root $\theta$, and let $T_1 = U_1 = S\alpha^4$ be the like function for the five roots; we have $T_1 = T + \theta^4$, or say $U_1 = U + \theta^4$. Write $-b_1$, $c_1$, $-d_1$, $e_1$, $-f_1$
% ----- printed page 280 (PDF 302) -----
for the symmetric functions of the five roots; $U_1$ will not involve $f_1$ and it will be the same function of $b_1, c_1, d_1, e_1$ that $U$ is of $b, c, d, e$, say we have
\[
U_1 = U(b_1, c_1, d_1, e_1).
\]
But we have
\[
b_1 = b - \theta,\ c_1 = c - b\theta,\ d_1 = d - c\theta,\ e_1 = e - d\theta;
\]
and thus the foregoing equation $U_1 = U + \theta^4$ becomes
\[
U(b - \theta,\ c - b\theta,\ d - c\theta,\ e - d\theta) = U + \theta^4;
\]
it is in fact easy to verify that, for the foregoing value of $U$, the terms in $\theta$, $\theta^2$, $\theta^3$ all vanish, and that the expression on the left-hand becomes $= U + \theta^4$. But attending only to the term in $\theta$, this is $= -\theta(\partial_b + b\partial_c + c\partial_d + d\partial_e) U$, $= -\theta \Delta U$; viz.\ this term vanishing we have $\Delta U = 0$, the result which was to be proved.

In the case of a unitary function, for instance $T = U = S\alpha^3\beta$, here introducing the new root $\theta$ we have $U_1 = U + \theta S\alpha^3 + \theta^3 S\alpha$; or there is here a term in $\theta$, and instead of $\Delta U = 0$, we have $\Delta U = S\alpha^3$, or the unitary function is not annihilated by $\Delta$.

The foregoing investigation is really quite general, and establishes the conclusion that $\Delta$ is an annihilator of every non-unitariant.

It is to be noticed that $\Theta_\sigma$ and $\Delta$ are operators, which leave each of them the degree unaltered but diminish the weight by unity: the operator $P - b\delta$, and another operator $\tfrac{1}{2}Q - b\omega$ which will be considered, increase each of them the degree by unity and also the weight by unity.

24.\ Coming now to the equation
\[
P - b\delta = S\sigma\partial_a,
\]
it is to be remarked that, if $\sigma' = \sigma$, the expression for $P$ ends in $q\partial_p$ where, as before, $q = a_\sigma$ is the highest coefficient in the operand; since the operand thus contains $q$, the next succeeding term in $r\partial_q$ would be not inoperative, and in order to include it in the expression of $P$ we may take $\sigma' = \sigma + 1$; we thus have
\[
P = b\partial_a + 2c\partial_b + 3d\partial_c + \ldots + (\sigma + 1) r\partial_q,
\]
or as we may more simply write it
\[
P = b\partial_a + c\partial_b + d\partial_c + \ldots + r\partial_q;
\]
the operation thus increases the extent by unity. The symbol $S\sigma\partial_a$ operating upon a symmetric function of the roots, gives, whatever may be the number of roots, the same symmetric function of the roots: and we see further that, operating upon a non-unitary function, it gives a non-unitary function of the roots. Hence $P - b\delta$ operating upon a non-unitariant gives a non-unitariant. I give an example.

25.\ Suppose, as before,
\[
T = U = S\alpha^4 = E = -4e + 4bd + 2c^2 - 4b^2c + b^4;
\]

% ----- printed page 281 (PDF 303) -----
here $\delta = 4$, and therefore
\[
P - 4\delta = b\partial_a + 2c\partial_b + 3d\partial_c + 4e\partial_d + 5f\partial_e - 4b.
\]
We have
\[
S\sigma\partial_a \cdot S\alpha^4 = 4S\alpha^5 = 4(-5f + 5be + 5cd - 5b^2d - 5bc^2 + 5b^3c - b^5),
\]
and this should therefore be the result of the operation $P - b\delta$: the calculation is

\begin{center}
\begin{tabular}{l|rrrrrrr}
 & $f$ & $be$ & $cd$ & $b^2d$ & $bc^2$ & $b^3c$ & $b^5$ \\
\hline
$b \cdot -12c + 8bd + 4c^2 - 4b^2c$ & & & & & & & \\
$2c \cdot 4d - 8bc + 4b^3$ & & & & & & & \\
$3d \cdot 4c - 4b^2$ & & & & & & & \\
$4e \cdot -4b$ & & & & & & & \\
$5f \cdot -4$ & & & & & & & \\
$-4b \cdot -4e + 4bd + 2c^2 - 4b^2c + b^4$ & & & & & & & \\
\hline
$4S\alpha^5$ & $-20$ & $+20$ & $+20$ & $-20$ & $-20$ & $+20$ & $-4$ \\
\end{tabular}
\end{center}

which is the right result.

We have seen that every non-unitariant is annihilated by $\Delta$; it at once appears that conversely every function of the coefficients which is annihilated by $\Delta$ is a non-unitariant: it is, in fact, a symmetric function of the roots, and unless it were a non-unitary function of the roots it would not be annihilated by $\Delta$. Non-unitariants are analogous to seminvariants; the precise relation between them will be shown further on.

26.\ We can, by an investigation similar to that for seminvariants, show that $P - b\delta$ operating upon a non-unitariant gives a non-unitariant. In fact, considering the two operations $\Delta$ and $P - b\delta$, we have
\[
\Delta(P - b\delta)^\dagger = \Delta(P - b\delta) + \Delta \cdot (P - b\delta),
\]
the meaning being that, if upon any operand $U$ we perform first the operation $P - b\delta$ and then the operation $\Delta$, this is equivalent to operating on $U$ with the sum of the two operations $\Delta(P - b\delta)$, and $\Delta \cdot P - b\delta$, the first of these symbols denoting the mere algebraical product of $\Delta$ and $P - b\delta$, the second of them the result of the operation $\Delta$ performed upon $P - b\delta$. We have similarly
\[
(P - b\delta)\Delta^\dagger = (P - b\delta)\Delta + (P - b\delta)\cdot\Delta.
\]

Hence observing that $\Delta(P - b\delta)$ and $(P - b\delta)\Delta$ are equal to each other, and subtracting, we have
\[
\Delta(P - b\delta)^\dagger - (P - b\delta)\Delta^\dagger = \Delta\cdot(P - b\delta) - (P - b\delta)\cdot\Delta.
\]
But from the values
\[
\Delta = \partial_b + b\partial_c + c\partial_d + \ldots
\]
and
\[
P - b\delta = b\partial_a + c\partial_b + d\partial_c + \ldots - b\delta,
\]
% ----- printed page 282 (PDF 304) -----
we find
\[
\Delta\cdot(P - b\delta) = a\partial_a + 2b\partial_b + 3c\partial_c + \ldots - b,
\]
\[
(P - b\delta)\cdot\Delta = b\partial_b + 2c\partial_c + \ldots,
\]
and thence
\[
\Delta\cdot(P - b\delta) - (P - b\delta)\cdot\Delta = a\partial_a + b\partial_b + c\partial_c + \ldots - \delta = 0,
\]
since $\delta$ is the degree in the coefficients. Hence writing down the operand $U$,
\[
\Delta\cdot(P - b\delta) U - (P - b\delta)\cdot\Delta U = 0,
\]
where for greater clearness I have inserted the dots, to show that $\Delta$ operates on $(P - b\delta) U$, and $(P - b\delta)$ on $\Delta U$. Taking $U$ to be a non-unitariant, we have $\Delta U = 0$; and this being so, the equation gives $\Delta\cdot(P - b\delta) U = 0$, viz.\ this shows that $(P - b\delta) U$ is a non-unitariant.

27.\ There is another symbol $\tfrac{1}{2}Q - b\omega$, which is precisely analogous to $P - b\delta$, viz.\ operating upon a non-unitariant, it gives a non-unitariant: $\omega$ is, as before, the weight of the function operated upon, and the expression of $Q$ is
\[
\tfrac{1}{2}Q = c\partial_b + 3d\partial_c + 6e\partial_d + \ldots + \tfrac{1}{2}\sigma(\sigma - 1)p\partial_{q-1},
\]
or say
\[
\tfrac{1}{2}Q = c\partial_b + 3d\partial_c + 6e\partial_d + \ldots.
\]
The proof is exactly similar, viz.\ we have to show that
\[
\Delta\cdot(\tfrac{1}{2}Q - b\omega) - (\tfrac{1}{2}Q - b\omega)\cdot\Delta = 0.
\]
We have
\[
\Delta\cdot(\tfrac{1}{2}Q - b\omega) = b\partial_b + 3c\partial_c + 6d\partial_d + \ldots - \omega,
\]
\[
(\tfrac{1}{2}Q - b\omega)\cdot\Delta = c\partial_c + 3d\partial_d + \ldots,
\]
and the difference of the two expressions is
\[
b\partial_b + 2c\partial_c + 3d\partial_d + \ldots - \omega = 0,
\]
since $\omega$ is the weight of the function operated upon. Hence, as before, if $U$ be a non-unitariant and therefore $\Delta U = 0$, we have $\Delta\cdot(\tfrac{1}{2}Q - b\omega) U = 0$, that is, $(\tfrac{1}{2}Q - b\omega) U$ is also a non-unitariant.

28.\ The symbol $\tfrac{1}{2}Q - b\omega$ has no simple expression in terms of $\partial_\alpha$, $\partial_\beta$, $\partial_\gamma$, \ldots, and the form varies with the number of the roots: thus for 3 roots, it is
\[
-\left\{\!\left(\frac{c\alpha^2 + 3d\alpha}{\alpha - \beta\cdot\alpha - \gamma} + b\alpha\!\right) \partial_\alpha + \&c.\!\right\},
\]
for 4 roots it is
\[
-\left\{\!\left(\frac{c\alpha^3 + 3d\alpha^2 + 6e\alpha}{\alpha - \beta\cdot\alpha - \gamma\cdot\alpha - \delta} + b\alpha\!\right)\partial_\alpha + \&c.\!\right\},
\]
for 5 roots it is
\[
-\left\{\!\left(\frac{c\alpha^4 + 3d\alpha^3 + 6e\alpha^2 + 10f\alpha}{\alpha - \beta\cdot\alpha - \gamma\cdot\alpha - \delta\cdot\alpha - \epsilon} + b\alpha\!\right)\partial_\alpha + \&c.\!\right\},
\]
and so on. It is not easy to find the effect of such a symbol upon a given symmetric function of the roots, nor in particular when the function is non-unitary is it easy to show generally that the result is non-unitary.

% ----- printed page 283 (PDF 305) -----
It is to be remarked that, if the function operated upon is of the degree $\delta$ in the roots, then we must for $\tfrac{1}{2}Q - b\omega$ take the expression with $\delta + 1$ roots; for instance, if the function be of the degree 5 in the roots, then, qua function of the coefficients, this contains $f$, and it must be operated on with
\[
\tfrac{1}{2}Q - b\omega = c\partial_b + 3d\partial_c + 6e\partial_d + 10f\partial_e + 15g\partial_f - b\omega,
\]
viz.\ this expression, as containing $g$, gives the 6-root expression for $\tfrac{1}{2}Q - b\omega$.

29.\ Suppose for instance the function operated upon is $F = S\alpha^5$; here taking the 6-root expression, this gives
\[
-5\left\{\!\left(\frac{c\alpha^5 + 3d\alpha^4 + 6e\alpha^3 + 10f\alpha^2 + 15g\alpha}{\alpha - \beta\cdot\alpha - \gamma\cdot\alpha - \delta\cdot\alpha - \epsilon\cdot\alpha - \zeta} + b\alpha\!\right) + \&c.\!\right\},
\]
or omitting for the moment the outside factor $-5$, the expression in $\{\ \}$ is easily seen to be
\[
= cH_4 + 3dH_3 + 6eH_2 + 10fH_1 + 15g\cdot 1 + bS\alpha^5,
\]
where $H_4$, $H_3$, $H_2$, $H_1$ denote the homogeneous functions of the degrees 4, 3, 2, 1 respectively: the values of these are obtained by adding together all the lines of the Table IV ($b$), all the lines of the Table III ($b$), \&c.: the terms exclusive of $bS\alpha^5$ thus are
\begin{align*}
& c(-e + 2bd + c^2 - 3b^2c + b^4) \\
& {} + 3d(-d + 2bc - b^3) \\
& {} + 6e(-c + b^2) \\
& {} + 10f(-b) \\
& {} + 15g\cdot 1
\end{align*}
and these are $= S\alpha^5\beta + S\alpha^4\beta^2 + S\alpha^3\beta^3$, as appears by the following calculation:

\begin{center}
\begin{tabular}{l|rrrr}
 & $S\alpha^5\beta$ & $S\alpha^4\beta^2$ & $S\alpha^3\beta^3$ & sum \\
\hline
$g$ & $+15$ & & & $+15$ \\
$bf$ & $-10$ & $-1$ & & $-11$ \\
$ce$ & $-7$ & $-6$ & $+2$ & $-?$ \\
$d^2$ & $+9$ & $+15$ & $-3$ & $-3$ \\
$b^2e$ & $-3$ & $-3$ & & \\
$bcd$ & $+1$ & $+1$ & $+1$ & \\
$b^4c$ & $+1$ & & & $+1$ \\
$b^2c^2$ & $+1$ & $+1$ & & \\
$c^3$ & & & & \\
$b^6$ & & & & \\
\end{tabular}
\end{center}

(Numerical entries in the columns of the original tabulation are partially illegible at scan resolution; the algebraic identity asserted by Cayley is $cH_4 + 3dH_3 + 6eH_2 + 10fH_1 + 15g = S\alpha^5\beta + S\alpha^4\beta^2 + S\alpha^3\beta^3$.)

% ----- printed page 284 (PDF 306) -----
The omitted term $bS\alpha^5$, that is, $-S\alpha\cdot S\alpha^5$, is $-S\alpha^6 - S\alpha^5\beta$; the addition hereof destroys therefore the non-unitary term $S\alpha^5\beta$, and thus the required expression, restoring the omitted factor $-5$, is $-5(-S\alpha^6 + S\alpha^4\beta^2 + S\alpha^3\beta^3)$, or say $= 5G - 5CE - 5D^2$, a non-unitary form: this then should be the result of the operation
\[
\tfrac{1}{2}Q - b\omega = c\partial_b + 3d\partial_c + 6e\partial_d + 10f\partial_e + 15g\partial_f - 5b,
\]
performed upon
\[
S\alpha^5 = F = -5f + 5be + 5cd - 5b^2d - 5bc^2 + 5b^3c - b^5.
\]

Performing the calculation so as to omit on each side a factor 5, it is to be shown that $G - CE - D^2$ is
\begin{align*}
& = c(e - 2bd - c^2 + 3b^2c - b^4) \\
& {} + 3d(d - 2bc + b^3) \\
& {} + 6e(c - b^2) \\
& {} + 10f(-1)+ 15g\cdot 1 \\
& {} - b(-f + be + cd - b^2d - bc^2 + b^3c - \tfrac{1}{5}b^5).
\end{align*}
Collecting the terms, and comparing the result with the expression for $G - CE - D^2$, we have:

\begin{center}
\begin{tabular}{l|rrrr}
 & RHS & $G$ & $-CE$ & $-D^2$ \\
\hline
$g$ & $-15$ & $-15$ & $-6$ & $-6$ \\
$bf$ & $+10$ & $+5$ & $+6$ & $+6$ \\
$ce$ & $-6$ & $-2$ & $+11$ & $-7$ \\
$d^2$ & $+3$ & $+3$ & $+3$ & $-3$ \\
$b^2e$ & $-5$ & $-13$ & $-12$ & $-4$ \\
$b cd$ & $-?$ & $+8$ & $+6$ & $+2$ \\
$b^2 c^2$ & $-?$ & $+9$ & $+2$ & $-1$ \\
$b^4 c$ & $-?$ & $-6$ & $-?$ & $-?$ \\
$b^6$ & $+1$ & $+1$ & & \\
\end{tabular}
\end{center}

(Several numerical entries in the original tabulation are too damaged for reliable reading; the identity $G - CE - D^2$ as asserted by Cayley is preserved as the column sum.)

The two expressions are thus identical.

30.\ Suppose again, 6 roots as before, and that the function operated upon is $S\alpha^3\beta^2$; we find $\partial_\alpha S\alpha^3\beta^2 = 3\alpha^2 S\beta^2 + 2\alpha S\beta^3 - 5\alpha^4$, and the general term is
\[
-\left\{\left(\frac{c\alpha^5 + 3d\alpha^4 + 6e\alpha^3 + 10f\alpha^2 + 15g\alpha}{\alpha - \beta\cdot\alpha - \gamma\cdot\alpha - \delta\cdot\alpha - \epsilon\cdot\alpha - \zeta} + b\alpha\right)(3\alpha^2 S\beta^2 + 2\alpha S\beta^3 - 5\alpha^4)\right\}.
\]

% ----- printed page 285 (PDF 307) -----
This gives
\begin{align*}
& -3(CH_2 + 3dH_1 + 6e + bS\alpha^3) S\alpha^2 \\
& {} - 2(CH_3 + 3dH_2 + 6eH_1 + 10f + bS\alpha^4) S\alpha^3 \\
& {} + 5(CH_4 + 3dH_3 + 6eH_2 + 10fH_1 + 15g + bS\alpha^5),
\end{align*}
which is found to be
\begin{align*}
& = -3(BD + C^2 + bS\alpha^3) S\alpha^2 \\
& {} - 2(BC + bS\alpha^4) S\alpha^3 \\
& {} + 5(BF + CE + D^2 + bS\alpha^5).
\end{align*}
Here
\[
bS\alpha^3 = -S\alpha\cdot S\alpha^3 = -S\alpha^4 - S\alpha^3\beta = -E - BD,
\]
\[
bS\alpha^4 = -S\alpha\cdot S\alpha^4 = -S\alpha^5 - S\alpha^4\beta = -D - BC',
\]
and
\[
bS\alpha^5 = -S\alpha\cdot S\alpha^5 = -S\alpha^6 - S\alpha^5\beta = -G - BF;
\]
the expression thus is
\begin{align*}
& = -3(-E + C^2)\cdot C \\
& \qquad -2(-D)\cdot D \\
& \qquad +5(-G + CE + D^2),\quad \text{that is,}\quad -3(-S\alpha^4 + S\alpha^2\beta^2)\,S\alpha^2 \\
&\hspace{4cm}-2(-S\alpha^3)\,S\alpha^3 \\
&\hspace{4cm}+5(-S\alpha^6 + S\alpha^4\beta^2 + S\alpha^3\beta^3).
\end{align*}
Here
\[
S\alpha^4\cdot S\alpha^2 = S\alpha^6 + 2S\alpha^4\beta^2,\quad = G + 2CE,
\]
\[
S\alpha^3\cdot S\alpha^3 = S\alpha^6 + 2S\alpha^3\beta^3,\quad = G + 2D^2,
\]
\[
S\alpha^2\beta^2\cdot S\alpha^2 = S\alpha^4\beta^2 + 3S\alpha^2\beta^2\gamma^2,\quad = CE + 3C^3;
\]
and the whole is
\begin{align*}
& = -3\{-G - CE + (CE + 3C^3)\} \\
& \qquad -2\{-G - 2D^2\} \\
& \qquad + 5(-G + CE + D^2),
\end{align*}
which is $= 5CE + 9D^2 - 9C^3$ (a non-unitary form). This then should be the value of
\[
\tfrac{1}{2}Q - b\omega = c\partial_b + 3d\partial_c + 6e\partial_d + 10f\partial_e + 15g\partial_f - 5b,
\]
operating upon
\[
S\alpha^3\beta^2 = CD = 5f - 5be + cd + b^2d - bc^2.
\]

31.\ There is for non-unitariants a theorem which is a much more simple form than the transformation of it afterwards obtained for seminvariants: viz.\ for any non-unitariant we have $\Delta U = 0 = (\partial_b + b\partial_c + c\partial_d + \ldots) U$; attending only to the portion $U'$ of $U$ which is of the highest degree, it is clear that we have $(b\partial_c + c\partial_d + \ldots) U' = 0$, and if we herein diminish the letters, then $(\partial_b + b\partial_c + \ldots) U'' = 0$, where $U''$ is what $U'$ becomes by a diminution of the letters; that is, $U''$ is a non-unitariant, viz.\ in any seminvariant, the terms of highest degree $U'$ are obtained from a non-unitariant $U''$ by a mere augmentation of the letters: e.g.\ $2e - 2bd + c^2$ is a non-unitariant weight 4; augmenting the letters, we have $2bf - 2ce + d^2$ which with a change of sign is the portion of highest degree of the non-unitariant $2g - 2bf + 2ce - d^2$.

% ----- printed page 286 (PDF 308) -----
\begin{center}
\textit{The MacMahon Form of Equation.} Art.\ Nos.\ 32 to 34.
\end{center}

32.\ The equation connecting the coefficients and the roots is here taken to be
\[
1 + \frac{b}{1}x + \frac{c}{1\cdot 2}x^2 + \frac{d}{1\cdot 2\cdot 3}x^3 + \ldots = \overline{1 - \alpha x}\cdot\overline{1 - \beta x}\cdot\overline{1 - \gamma x}\cdots
\]
As to this it may be remarked that, if we had started with a form of the $n$th order with binomial coefficients,
\[
1 + \frac{n}{1}bx + \frac{n\cdot n-1}{1\cdot 2}cx^2 + \frac{n\cdot n-1\cdot n-2}{1\cdot 2\cdot 3}dx^3 + \ldots = \overline{1 - \alpha x}\cdot\overline{1 - \beta x}\cdot\overline{1 - \gamma x}\cdots\ (n \text{ factors}),
\]
then writing herein $\dfrac{x}{n}$ for $x$, and also $n\alpha$, $n\beta$, $n\gamma$, \ldots, for $\alpha$, $\beta$, $\gamma$, \ldots, and putting ultimately $n = \infty$, we have the form in question.

We pass from the ordinary form to the MacMahon form, by writing for
\[
b,\ c,\ d,\ e,\ \ldots\quad \frac{b}{1},\ \frac{c}{1\cdot 2},\ \frac{d}{1\cdot 2\cdot 3},\ \frac{e}{1\cdot 2\cdot 3\cdot 4},\ \ldots\quad\text{or say}\quad b,\ \frac{c}{2},\ \frac{d}{6},\ \frac{e}{24},\ \frac{f}{120},\ \frac{g}{720},\ \ldots
\]

All the results obtained for the ordinary form will, after making therein this change, apply to the new form. We thus find
\[
\Theta_\sigma = \sigma c\partial_b + (\sigma - 1)\, 2b\partial_c + (\sigma - 2)\, 3c\partial_d + \ldots + \quad \sigma p\partial_q,
\]
\[
\Theta_{\sigma'} = \sigma' c\partial_b + (\sigma' - 1)\, 2b\partial_c + (\sigma' - 2)\, 3c\partial_d + \ldots + (\sigma' - \sigma + 1)\, \sigma p\partial_q,
\]
where
\[
\Theta_{\sigma'} - \Theta_\sigma = (\sigma' - \sigma)\Delta,
\]
or say
\[
\Delta = \partial_b + 2b\partial_c + 3c\partial_d + \ldots + \sigma p\partial_q,
\]
or say
\[
\Delta = \partial_b + 2b\partial_c + 3c\partial_d + \ldots.
\]
Also
\[
P = b\partial_a + c\partial_b + d\partial_c + \ldots + r\partial_q,
\]
or say
\[
P = b\partial_a + c\partial_b + d\partial_c + \ldots,
\]
\[
Q = c\partial_a + 2d\partial_b + 6e\partial_c + \ldots,
\]
or say
\[
\tfrac{1}{2}Q = c\partial_b + 2d\partial_c + \ldots.
\]

The change $\alpha$, $\beta$, $\gamma$, \ldots\ into $n\alpha$, $n\beta$, $n\gamma$, \ldots\ would change $S\partial_a$, $S\alpha\partial_a$, $S\alpha^2\partial_a$ into $n S\partial_a$, $S\alpha\partial_a$, $n S\alpha^2\partial_a$ respectively ($n = \infty$): but this change is, in fact, compensated for by the introduction into the formulae of the binomial coefficients as above; it is $-S\alpha$, $S\alpha\beta$, \ldots\ not $-nS\alpha$, $n^2 S\alpha\beta$, \ldots\ which are equal to $b$, $\tfrac{1}{2}c$, \ldots; and the conclusion is that we have to retain without alteration the symbols $S\partial_a$, $S\alpha\partial_a$, $S\alpha^2\partial_a$: thus in the new form as in the old one, we have $\Theta_4 S\alpha^4 = -S\partial_a\cdot S\alpha^4 = -4S\alpha^3$, see the example ante No.\ 23.

33.\ In the new form, a non-unitariant is annihilated by the operator
\[
\Delta,\ = \partial_b + 2b\partial_c + 3c\partial_d + \ldots,
\]
% ----- printed page 287 (PDF 309) -----
and conversely any function annihilated by $\Delta$ is a non-unitariant; comparing here with the subsequent theory of seminvariants, this is in fact the theorem that a non-unitariant is the same thing as a seminvariant; or to state this more explicitly: for the MacMahon form of equation, a function of the coefficients which is a non-unitary symmetric function of the roots is a seminvariant.

I consider for instance the Table VI ($b$), but attend only to the non-unitary portions thereof, viz.\ the lines $G$, $CE$, $D^2$, $C^3$: I convert these into columns, at the same time changing the arrangement of the headings $g$, $bf$, $ce$, \&c., from CO to AO: and then making the foregoing change $b, c, d, e, f, g$ into $b,\ \tfrac{c}{2},\ \tfrac{d}{6},\ \tfrac{e}{24},\ \tfrac{f}{120},\ \tfrac{g}{720}$, but to avoid fractions multiplying the whole by 720, I form the table

\begin{center}
\begin{tabular}{l|rrrr}
 & $C^3$ & $D^2$ & $CE$ & $G$ \\
\hline
$1\cdot g$ & $-2$ & $-3$ & $+6$ & $-6$ \\
$6\cdot bf$ & $+2$ & $-3$ & $-6$ & $+6$ \\
$15\cdot ce$ & $-2$ & $-3$ & $+2$ & $+6$ \\
$20\cdot d^2$ & $+1$ & $+3$ & $-3$ & $+3$ \\
$30\cdot b^2 e$ & & $-3$ & $-2$ & $-6$ \\
$60\cdot bcd$ & & $-3$ & $+4$ & $-12$ \\
$90\cdot c^3$ & & $+1$ & $-2$ & $-2$ \\
$120\cdot b^3 d$ & & & $-2$ & $+6$ \\
$180\cdot b^2 c^2$ & & & $+1$ & $+9$ \\
$360\cdot b^4 c$ & & & & $-6$ \\
$720\cdot b^6$ & & & & $+1$ \\
\hline
& $[d^2]$ & $[c^3]$ & $[b^2 c^2]$ & $[b^6]$ \\
\end{tabular}
\end{center}

which is to be read according to the columns: and observe that the outside left-hand numbers are to be multiplied into the numbers of each column: thus the first column is to be read
\[
C^3 = S\alpha^2\beta^2\gamma^2 = \tfrac{1}{720}(-2g + 12bf - 30ce + 20d^2),
\]
the second column is to be read
\[
D^2 = S\alpha^3\beta^3 = \tfrac{1}{720}(3g - 18bf - \ldots + 90c^3),
\]
and so on.

By what precedes, the columns are seminvariants,---as afterwards explained, ``blunt'' seminvariants; and they are named as such by the outside bottom line of symbols with a $[\ ]$; viz.
\[
[d^2] = (-2g + 12bf - 30ce + 20d^2),\quad [c^3] = (3g - 18bf - \ldots + 90c^3),\ \&c.,^*
\]
% ----- printed page 288 (PDF 310) -----
where it will be observed that the symbol within the $[\ ]$ is, in fact, the power-ender which is in AO the lowest term of the column; and further that this is also the conjugate of the capital letter symbol at the head of the column.

The ($b$) Tables I to X, with only the change $b, c, d, e, \ldots$ into $b,\ \tfrac{c}{2},\ \tfrac{d}{6},\ \tfrac{e}{24},\ \ldots$ are given in my paper, ``Tables of the Symmetric Functions of the Roots to the degree 10, for the form
\[
1 + bx + \frac{c x^2}{1\cdot 2} + \ldots = (1 - \alpha x)(1 - \beta x)(1 - \gamma x)\ldots,
\]
\textit{American Mathematical Journal}, t.\ \textsc{vii}.\ (1885), pp.\ 47--56, [829].''

34.\ By what precedes, it appears that $P - b\delta$ operating on a seminvariant gives a seminvariant, and that $Q - 2b\omega$ operating on a seminvariant gives a seminvariant: these operators will be further considered in the development of the theory of seminvariants. We see further that $\tfrac{1}{2}\Delta$, $= b\partial_c + 3c\partial_d + 6d\partial_e + \ldots$, operating on a seminvariant gives sometimes but not always a seminvariant, e.g.
\[
(b\partial_c + 3c\partial_d + 6d\partial_e)(e - 4bd - 3c^2 + 12b^2c - 6b^4) = 6(d - 3bc + 2b^3).
\]

\begin{center}
\textit{Seminvariants---the I-and-F Problem, and Solution by Square Diagrams.}\\
Art.\ Nos.\ 35 to 47.
\end{center}

35.\ Writing
\begin{align*}
1 &= 1, \\
b_1 &= b + \theta, \\
c_1 &= c + 2b\theta + \theta^2, \\
d_1 &= d + 3c\theta + 3b\theta^2 + \theta^3, \\
e_1 &= e + 4d\theta + 6c\theta^2 + 4b\theta^3 + \theta^4, \\
&\;\&c.,
\end{align*}
then there are functions of the unsuffixed letters which remain unaltered if for these we substitute the suffixed letters: any such function is termed a seminvariant. We have for instance
\[
c_1 = c + 2b\theta + \theta^2,\qquad\text{i.e.,}\quad c_1 - b_1^2 = c - b^2,
\]
\[
-b_1^2 = -b^2 - 2b\theta - \theta^2,
\]
\[
d_1 = d + 3c\theta + 3b\theta^2 + \theta^3,\qquad d_1 - 3b_1 c_1 + 2b_1^3 = d - 3bc + 2b^3,
\]
\[
-3b_1 c_1 = -3bc - 3c\theta - 6b^2\theta - 6b\theta^2 - 3\theta^3,
\]
\[
+ 2b_1^3 = 2b^3 + 6b^2\theta + 6b\theta^2 + 2\theta^3,
\]
% ----- printed page 289 (PDF 311) -----
and thus $c - b^2$, $d - 3bc + 2b^3$ are seminvariants; they are, in fact, the first and second terms of the series
\begin{align*}
& c - b^2, \\
& d - 3bc + 2b^3, \\
& e - 4bd + 6b^2 c - 3b^4, \\
& f - 5be + 10b^2 d - 10b^3 c + 4b^5, \\
& g - 6bf + 15ce - 20b^2 e + 15b^4 c - 5b^6, \\
& \;\vdots
\end{align*}
where the law is obvious; the numbers in each line are binomial coefficients except the last number, which is the next binomial coefficient diminished by unity. The successive terms are, in fact, what $c_1$, $d_1$, $e_1$, $f_1$, $g_1$, \ldots\ become upon writing therein $\theta = -b$.

36.\ Any rational and integral function of these forms is a seminvariant, and it is to be observed that we can form functions for which (by the destruction of terms of a higher degree) there is a diminution of degree; for instance,
\[
(e - 4bd + 6b^2 c - 3b^4) + 3(c - b^2)^2
\]
gives a seminvariant $e - 4bd + 3c^2$.

It is important to remark that a seminvariant is completely determined by its non-unitary terms; thus for $e - 4bd + 3c^2$, the non-unitary terms are $e + 3c^2$, and for this writing $e_1 + 3c_1^2$, and for $e_1$, $c_1$ substituting their above values for $\theta = -b$, we reproduce the original value $e - 4bd + 3c^2$.

37.\ It is at once seen that a seminvariant is reduced to zero by the operation $\Delta,\ = \partial_b + 2b\partial_c + 3c\partial_d + \ldots$, or say that $\Delta$ is an annihilator of a seminvariant; in fact, if in any function of $b, c, d, \ldots$ we write for these the suffixed letters $b_1, c_1, d_1, \ldots$, then the coefficient of $\theta$ herein is at once found by operating on the function of $(b, c, d, \ldots)$ with $\Delta$, and therefore in the case of a seminvariant the result of this operation must be $= 0$. And conversely, every function of $(b, c, d, \ldots)$ which is reduced to zero by the operation $\Delta$ is a seminvariant.

38.\ For a given weight, the number of seminvariants is equal to the excess of the number of terms of that weight above the number of terms of the next preceding weight, or what is the same thing, it is equal to the number of power-enders of the given weight. More definitely, considering the terms of a seminvariant as arranged in AO, we have seminvariants the finals whereof are the several power-enders of the given weight: and we arrange the seminvariants inter se by taking these power-enders in AO: thus for the weight 6, we have seminvariants $[d^2]$, $[c^3]$, $[b^2 c^2]$, $[b^6]$ ending in these terms respectively. We may, if we please, consider all these seminvariants as beginning with $g$, or say the forms may be taken to be
% ----- printed page 290 (PDF 312) -----
$g(\text{ao}) d^2,\ g(\text{ao}) c^3,\ g(\text{ao}) b^2 c^2,\ g(\text{ao}) b^6.$ Such forms are, in fact, furnished by the MacMahon equation: viz.\ up to the weight 6, we thus have for the present purpose

\begin{center}
\textit{[Tabulation of $[c]$, $[d]$, $[e]$, $[f]$, $[c^2]$, $[b^2]$ \&c.\ extending columns through weight 6; the structure follows the format described in the text; numerical entries are tabulated in the Cayley original.]}
\end{center}

\medskip

\begin{center}
\begin{tabular}{l|rrrr}
 & $[d^2]$ & $[c^3]$ & $[b^2 c^2]$ & $[b^6]$ \\
\hline
$1\cdot g$ & $-2$ & $+3$ & $-6$ & $-6$ \\
$6\cdot bf$ & $+12$ & $-18$ & $+6$ & $+6$ \\
$15\cdot ce$ & $-30$ & $-45$ & $+2$ & $-3$ \\
$20\cdot d^2$ & $+20$ & $+60$ & $-3$ & $+3$ \\
$30\cdot b^2 e$ & & & $-2$ & $-6$ \\
$60\cdot bcd$ & & & $+4$ & $-12$ \\
$90\cdot c^3$ & & $+90$ & $-2$ & $-2$ \\
$120\cdot b^3 d$ & & & $-2$ & $+6$ \\
$180\cdot b^2 c^2$ & & & $+1$ & $+9$ \\
$360\cdot b^4 c$ & & & & $+1$ \\
$720\cdot b^6$ & & & & $+1$ \\
\end{tabular}
\end{center}

\noindent read for instance
\[
[d^2] = (-2g + 12bf - 30ce + 20d^2),
\]
\[
[c^3] = (3g - 18bf - 45ce + 60d^2 + 90c^3),
\]
\&c.

I say that $[d^2]$, $[c^3]$, $[b^2 c^2]$, $[b^6]$ are ``specific'' when they are regarded as standing for these tabulated functions; but in general I take them to be ``indefinite,'' that is, I regard them as denoting (as above) any seminvariants ending in $d^2$, $c^3$, $b^2 c^2$, $b^6$ respectively.

% ----- printed page 291 (PDF 313) -----
39.\ The seminvariant $[d^2]$ is of the form $(g\,\mathrm{ao}\,d^2)$, including those terms which are in CO not superior to $g$ and in AO not inferior to $d^2$: by a combination of $[d^2]$ and $[c^3]$, we obtain a seminvariant $(ce\,\mathrm{ao}\,c^3)$ containing terms which are in CO not superior to $ce$ and in AO not inferior to $c^3$: similarly, from $[d^2]$, $[c^3]$, $[b^2 c^2]$ we obtain a seminvariant $(d^2\,\mathrm{ao}\,b^2 c^2)$; and from the four forms a seminvariant $(c^3\,\mathrm{ao}\,b^6)$: these four seminvariants

\begin{center}
\begin{tabular}{l|rrrr}
 & $(g\,\mathrm{ao}\,d^2)$ & $(ce\,\mathrm{ao}\,c^3)$ & $(d^2\,\mathrm{ao}\,b^2 c^2)$ & $(c^3\,\mathrm{ao}\,b^6)$ \\
\hline
$g$ & $+1$ & & & \\
$bf$ & $-6$ & & & \\
$ce$ & $+15$ & $+1$ & & \\
$d^2$ & $-10$ & $-1$ & $+1$ & \\
$b^2 e$ & & $+2$ & $-3$ & \\
$bcd$ & & $+4$ & $-3$ & \\
$c^3$ & & & $+1$ & $-1$ \\
$b^3 d$ & & & $+1$ & \\
$b^2 c^2$ & & & & $+3$ \\
\end{tabular}
\end{center}

are said to be ``sharp'' seminvariants: viz.\ considering the final as given, a sharp seminvariant is one having an initial which is in CO as low as possible; or considering the initial as given, it is one having a final which is in AO as high as possible. A seminvariant which is not sharp is said to be ``blunt.''

40.\ The sharp seminvariants are in general designated as above, $(g\,\mathrm{ao}\,d^2)$, \&c.: but it is sometimes convenient to give the numerical coefficients of the initial and final terms respectively: as to this, it is to be noticed that the coefficient of the initial term is in most cases, but not always, $= 1$,---we might of course take it to be always $= 1$, but we should then in the excepted cases have fractional coefficients, and it is better to avoid this by giving a proper value to the numerical coefficient of the initial term; the numerical coefficient of the final term is in general different from $+1$, and it is not in general a multiple of the numerical coefficient of the initial term. As an instance take $dh\,\mathrm{ao}\,b^2 e^2$, the more complete expression of which is $4dh\,\mathrm{ao}\,-35 b^2 e^2$. The sharp seminvariants up to the weight 12 are designated in this more complete form in the table post No.\ 62.

41.\ In the calculation of the sharp seminvariants by elimination as above, it will be noticed how unitary terms disappear: thus in combining $[d^2]$ and $[c^3]$ so as to get rid of $g$, the term $bf$ disappears of itself, and we have as above the form

% ----- printed page 292 (PDF 314) -----
$(ce\,\mathrm{ao}\,c^3)$ beginning with the non-unitary term $ce$. We may, in fact, write $b = 0$; we thus have
\[
[d^2] = -2g - 30ce + 20d^2,
\]
\[
[c^3] = 3g - 45ce + 60d^2 + 90c^3,
\]
giving $3[d^2] + 2[c^3] = -180(ce - d^2 - c^3)$, and then $ce - d^2 - c^3$, putting therein for $c, d, e$ the values $c - b^2$, $d - 3bc + 2b^3$, $e - 4bd + 6b^2c - 3b^4$, gives the complete value \textit{ut supra}, $ce - d^2 - b^2 e + 2bcd - c^3$, and we thus see \textit{a priori} that this contains no term $bf$, but in fact begins with $ce$. And in carrying out this process for any higher given weight, it is proper also to arrange the non-unitary terms not in AO but in CO, and then in each case beginning with the terms highest in CO and eliminating as many as possible of these terms we obtain the sharp seminvariant. Consider for instance the weight 12: taking the finals in AO, we have here
\[
(m\,\mathrm{ao}\,g^2),\ (m\,\mathrm{ao}\,cf^2),\ (m\,\mathrm{ao}\,e^3),\ (m\,\mathrm{ao}\,b^2 f^2),\ldots
\]
the initials in CO are $m, ck, dj, ei, \ldots$ and it might at first sight appear that the foregoing process of elimination would lead to the forms $(m\,\mathrm{ao}\,g^2)$, $(ck\,\mathrm{ao}\,cf^2)$, $(dj\,\mathrm{ao}\,e^3)$, $(ei\,\mathrm{ao}\,b^2 f^2)$, \ldots; we in fact have the form $(m\,\mathrm{ao}\,g^2)$; and if from $(m\,\mathrm{ao}\,g^2)$ and $(m\,\mathrm{ao}\,cf^2)$ we eliminate $m$, we obtain the form $(ck\,\mathrm{ao}\,cf^2)$; but we cannot have a form $(dj\,\mathrm{ao}\,e^3)$ (for a form beginning with $dj$ is of necessity of the degree 4 at least); what happens is that when from $(m\,\mathrm{ao}\,g^2)$, $(m\,\mathrm{ao}\,cf^2)$ and $(m\,\mathrm{ao}\,e^3)$ we eliminate $m$ and $ck$, the next term $dj$ disappears of itself, and (the following term $ei$ not disappearing) the resulting form is $(ei\,\mathrm{ao}\,e^3)$: to obtain a form beginning with $dj$ we must use the fourth form $(m\,\mathrm{ao}\,b^2 f^2)$, and we thence obtain $(dj\,\mathrm{ao}\,b^2 f^2)$. Arranging the initials in CO and the finals in AO, we thus have

\begin{center}
\begin{tabular}{ccc}
$m$ &$\longrightarrow$& $g^2$ \\
$ck$ &$\longrightarrow$& $cf^2$ \\
$dj$ &$\longrightarrow$& $b^2 f^2$ \\
$ei$ &$\longrightarrow$& $e^3$ \\
\end{tabular}
\end{center}

that is, we have the sharp seminvariants $m\,\mathrm{ao}\,g^2$, $ck\,\mathrm{ao}\,cf^2$, $ei\,\mathrm{ao}\,e^3$, $dj\,\mathrm{ao}\,b^2 f^2$, \ldots; these are the results given by the MacMahon linkage as will be explained further on, but I will first approach the question from a different side.

42.\ It has been seen that we have $\Delta,\ = \partial_b + 2b\partial_c + 3c\partial_d + \ldots$, as the annihilator of a seminvariant. Considering in the first place the entire set of terms, say for the weight 6, $g(\text{ao}) b^6$, we assume for a seminvariant the sum of these each multiplied by an arbitrary coefficient; the number of coefficients is equal to the number of terms of $g(\text{ao}) b^6$. Operating with $\Delta$, we obtain a function of the next inferior weight 5, containing all the terms of $g(\text{ao}) b^6$, that is, of $f(\text{ao}) b^5$, each
% ----- printed page 293 (PDF 315) -----
term multiplied by a linear function (with mere numerical factors) of the arbitrary coefficients: the expression thus obtained must be identically $= 0$; and we thus find between the arbitrary coefficients a number of linear relations equal to the number of terms of $f(\text{ao}) b^5$: these relations are independent; or it is on the supposition that they are so, that the number of coefficients which remain arbitrary will be $11 - 7 = 4$, agreeing with the number of the seminvariants $[d^2]$, $[c^3]$, $[b^2 c^2]$, $[b^6]$; whereas if the relations were not independent, there would be a larger number of seminvariants.

But if, instead of the whole set $g(\text{ao}) b^6$, we consider a set $(g\,\mathrm{ao}\,d^2)$ or say $(ce\,\mathrm{ao}\,c^3)$ and assume for a seminvariant the sum of these terms each multiplied by an arbitrary coefficient, then operating as before with $\Delta$ we obtain between the arbitrary coefficients a number of relations equal to that of the terms $D(ce\,\mathrm{ao}\,c^3)$, and if this be less by unity than the number of the terms of $ce\,\mathrm{ao}\,c^3$, say if we have $(1 - D)(ce\,\mathrm{ao}\,c^3) = 1$, then there will be a single seminvariant $ce\,\mathrm{ao}\,c^3$. We, in fact, find $(1-D)$: $(g\,\mathrm{ao}\,d^2)$, $(ce\,\mathrm{ao}\,c^3)$, $(d^2\,\mathrm{ao}\,b^2 c^2)$, $(c^3\,\mathrm{ao}\,b^6)$, each $= 1$, and thus establish the existence of the foregoing seminvariants $g\,\mathrm{ao}\,d^2$, $ce\,\mathrm{ao}\,c^3$, $d^2\,\mathrm{ao}\,b^2 c^2$, $c^3\,\mathrm{ao}\,b^6$. And similarly if in any case we have $(1 - D)(I\,\mathrm{ao}\,F) = 2$ or any larger number, then we have 2 or more seminvariants $I\,\mathrm{ao}\,F$.

43.\ It will be convenient to write down at once the system of square diagrams for the several weights 2 to 16; each of these may theoretically be obtained by a direct process of calculation such as I exhibit for the weight 10, but the labour would be very great indeed, and I have in fact formed the squares for the weights 11 to 16, not in this manner but by the MacMahon linkage.

\textit{[Figure: Square diagrams for weights $w=2$ through $w=10$. Each diagram lists power-enders along the bottom edge (in AO) and non-unitary initials along the left edge (in CO), with $1$'s placed in cells where the corresponding seminvariant exists; the dexter diagonal carries $1$'s through $w=9$. For $w=10$ there are figures $1$ and $2$ in the dexter diagonal denoting multiplicities, with inversions among $(c^2 g, f^2)$ and $(cdf, ce^2)$. Original layout reproduced from scan; cf.\ p.\ 296 below.]}

% ----- printed page 294 (PDF 316) -----
\textit{[Figure: Square diagrams for $w=8$, $w=9$, $w=10$. The $w=8$ diagram has initials $i$, $cg$, $df$, $e^2$, $c^2 e$, $cd^2$, $c^4$; finals $b^8$, $b^6 c$, $b^4 c^2$, $b^2 c^3$, $c^4$, $b^4 d^2$, $b^2 d^2$. The $w=9$ diagram has initials $j$, $ch$, $dg$, $ef$, $c^2 f$, $cde$, $d^3$, $c^3 d$. The $w=10$ diagram has initials $k$, $ci$, $dh$, $eg$, $f^2$, $c^2 g$, $cdf$, $ce^2$, $d^2 e$, $c^3 e$, $c^2 d^2$, $c^4$ and finals $b^{10}, \ldots, c^5$; the entries $1$ and $2$ are placed along the dexter diagonal.]}

% ----- printed page 295 (PDF 317) -----
The subsequent squares $w = 11$ to 16 are, for convenience, given at the end of the present memoir (pp.\ 331 et seq.).

44.\ It is to be observed that in each square the outside left-hand terms are the non-unitaries in CO and the outside bottom terms are the power-enders in AO. I have inside each square written down only the significant numbers, but we might fill up the whole square. For instance, when $w = 7$, the filled-up square would be

\begin{center}
\begin{tabular}{c|cccc}
 & 1 & 2 & 3 & 4 \\
\hline
$h$ & 0 & 1 & 2 & 3 \\
$cf$ & 0 & 1 & 2 & 3 \\
$de$ & 0 & 0 & 1 & 2 \\
$c^2 d$ & 0 & 0 & 0 & 1 \\
\hline
 & $bd^2$ & $bc^3$ & $b^3 c^2$ & $b^7$ \\
\end{tabular}
\end{center}

where in the first column the numbers relate to the sets $h\,\mathrm{ao}\,bd^2$, $cf\,\mathrm{ao}\,bd^2$, $de\,\mathrm{ao}\,bd^2$ and $c^2 d\,\mathrm{ao}\,bd^2$ (this last set $c^2 d\,\mathrm{ao}\,bd^2$ is non-existent since $c^2 d$ is in AO inferior to $bd^2$, i.e.\ as well for the set as for the diminished set, number of terms is $= 0$, and we have for the compartment $0 - 0 = 0$). And similarly for the remaining three columns. The process of thus filling up the whole square is a direct and non-tentative one, and the conclusions to which the numbers lead are as follows: col.\ 1, the final being $bd^2$, the initial cannot be $c^2 d$, $de$ or $cf$, but taking it to be $h$, we have the seminvariant $h\,\mathrm{ao}\,bd^2$. Col.\ 2, the final being $bc^3$ the initial cannot be $c^2 d$ or $de$, but taking it to be $cf$ we have the seminvariant $cf\,\mathrm{ao}\,bc^3$: it may be added that the top number 2 shows that there are two seminvariants $h\,\mathrm{ao}\,bc^3$, these are of course the foregoing ones $h\,\mathrm{ao}\,bd^2$ and $cf\,\mathrm{ao}\,bc^3$. Similarly, col.\ 3, the final being $b^3 c^2$, the initial cannot be $c^2 d$, but taking it to be $de$, we have the seminvariant $de\,\mathrm{ao}\,b^3 c^2$, and col.\ 4, we have the seminvariant $c^2 d\,\mathrm{ao}\,b^7$.

For the several weights up to 9, we have simply units in the dexter diagonal of each square, viz.\ the non-unitaries in CO correspond to the power-enders in AO, or the sharp seminvariants are $c\,\mathrm{ao}\,b^2$, $d\,\mathrm{ao}\,b^3$, \&c. See post, Table of Reductions, No.\ 62, which exhibits these correspondences.

45.\ For the weight 10, we have deviations: the figures 1 and 2 denote as follows:

\begin{center}
\begin{tabular}{rcl}
$k$ & $\mathrm{ao}$ & $b^{10}$,\quad ``1'' \\
$ci$ & $\mathrm{ao}$ & $b^8 c$,\quad ``1'' \\
$dh$ & $\mathrm{ao}$ & $b^6 c^2$,\quad ``1'' \\
$eg$ & $\mathrm{ao}$ & $b^4 c^3$,\quad ``1'' \\
$f^2$ & $\mathrm{ao}$ & $b^4 d^2$,\quad ``1'' \\
$c^2 g$ & $\mathrm{ao}$ & $b^2 c d^2$,\quad ``2'' \\
$cdf$ & $\mathrm{ao}$ & $b^2 c d^2$,\quad ``2'' \\
$ce^2$ & $\mathrm{ao}$ & $b^2 c d^2$,\quad ``1'' \\
$d^2 e$ & $\mathrm{ao}$ & $c^2 d^2$,\quad ``1'' \\
$c^3 e$ & $\mathrm{ao}$ & $c^5$,\quad ``1'' \\
$c^2 d^2$ & $\mathrm{ao}$ & $b^6 c^2$,\quad ``1'' \\
$c^4$ & $\mathrm{ao}$ & $b^4 c^3$,\quad ``1'' \\
\end{tabular}
\end{center}

% ----- printed page 296 (PDF 318) -----
and they indicate the sharp seminvariants $k\,\mathrm{ao}\,f^2$, $ci\,\mathrm{ao}\,ce^2$, \&c.: where observe that the power-enders being in AO as before, the non-unitaries are not in CO, but we have inversions $(c^2 g, f^2)$ and $(cdf, ce^2)$.

In particular, $(1 - D)(f^2\,\mathrm{ao}\,c^2 d^2) = 1$ indicates the seminvariant $f^2\,\mathrm{ao}\,c^2 d^2$;
\[
(1 - D)(c^2 g\,\mathrm{ao}\,b^2 c d^2) = 2,
\]
means in the first instance that there are 2 seminvariants $c^2 g\,\mathrm{ao}\,b^2 c d^2$, but here the set $c^2 g\,\mathrm{ao}\,b^2 c d^2$ includes as part of itself the set $f^2\,\mathrm{ao}\,c^2 d^2$; so that, if $c^2 g\,\mathrm{ao}\,b^2 c d^2$ is used to denote any particular form, then the general form is $c^2 g\,\mathrm{ao}\,b^2 c d^2$ plus an arbitrary multiple of $f^2\,\mathrm{ao}\,c^2 d^2$, and we have thus virtually a single form $c^2 g\,\mathrm{ao}\,b^2 c d^2$. And similarly, the set $cdf\,\mathrm{ao}\,b^4 d^2$ includes as part of itself the set $ce^2\,\mathrm{ao}\,c^5$; and thus the general form $cdf\,\mathrm{ao}\,b^4 d^2$ is $=$ particular form plus an arbitrary multiple of $ce^2\,\mathrm{ao}\,c^5$, or we have virtually a single form $cdf\,\mathrm{ao}\,b^4 d^2$.

I remark that it would be allowable to take as a standard form of $c^2 g\,\mathrm{ao}\,b^2 c d^2$, a form not containing any term in $f^2$, and similarly for the standard form of $cdf\,\mathrm{ao}\,b^4 d^2$ a form not containing any term in $ce^2$; but this is not done in the tables.

46.\ The diagram for weight 10 is constructed by the following calculation; viz.\ in col.\ 1 we calculate $(1 - D)(k\,\mathrm{ao}\,f^2)$ and for this purpose write down the terms of $k\,\mathrm{ao}\,f^2$, and $D(k\,\mathrm{ao}\,f^2)$ in CO: in col.\ 2 we calculate $(1 - D)(ci\,\mathrm{ao}\,ce^2)$, and for this purpose write down the terms of $k\,\mathrm{ao}\,ce^2$ and $D(k\,\mathrm{ao}\,ce^2)$ in CO, the terms of $ci\,\mathrm{ao}\,ce^2$ and $D(ci\,\mathrm{ao}\,ce^2)$ being thence found by rejecting the terms $k$, $bj$ and the term $j$ at the head of the two halves of the column. So in col.\ 3 we calculate $(1 - D)(dh\,\mathrm{ao}\,b^2 e^2)$, and for this purpose write down the terms of $(k\,\mathrm{ao}\,b^2 e^2)$ and $D(k\,\mathrm{ao}\,b^2 e^2)$ in CO, and for $dh\,\mathrm{ao}\,b^2 e^2$ and $D(dh\,\mathrm{ao}\,b^2 e^2)$ reject the terms $k$, $bj$, $ci$, $b^2 i$ and $j$, $bi$ at the head of the two halves of the column. And so for the remaining columns. It is to be remarked that there is in each successive column a continually increasing number of terms to be rejected; by a properly devised variation of the algorithm it would have been possible to avoid writing down these terms at all, but for greater clearness I have inserted them.

\bigskip
\noindent\textit{[The original page here carries an extensive multi-column tabulation, headed by $k$, $bj$, $ci$, $dh$, $eg$, $f^2$, \&c., listing for each successive column ($w=10$) the terms of the set and of its $D$-image, with bottom marginal counts of the form ``11 -- 7 = 4'', ``14 -- 13 = 1'', \&c., showing the differences whose unit values yield the sharp seminvariants. The tabulation is too damaged at scan resolution for reliable cell-by-cell reproduction; the algorithm is fully described in the surrounding text.]}

% ----- printed page 298 (PDF 320) -----
47.\ As to the first of the foregoing inversions $c^2 g$, $f^2$, it is proper to remark, that filling up two compartments of the square we have

\begin{center}
\begin{tabular}{c|cc}
 & $c^2 d^2$ & $b^2 c d^2$ \\
\hline
$c^2 g$ & $(\imath)$ & \\
$f^2$ & & $(\imath)$ \\
\end{tabular}
\end{center}

where the meaning of the numbers $(\imath, \imath)$ has to be considered: the first $(\imath)$ seems to indicate a seminvariant $c^2 g\,\mathrm{ao}\,c^2 d^2$, but there is in fact no such form, what it really indicates is a form $0\cdot c^2 g + f^2\,\mathrm{ao}\,c^2 d^2$, that is, $f^2\,\mathrm{ao}\,c^2 d^2$; and similarly, the second $(\imath)$ seems to indicate a seminvariant $f^2\,\mathrm{ao}\,b^2 c d^2$, but there is in fact no such form, what it really indicates is $f^2\,\mathrm{ao}\,c^2 d^2 + 0\cdot b^2 c d^2$, that is, $f^2\,\mathrm{ao}\,c^2 d^2$. The explanation is correct, but to make it perfectly clear some further developments would be required. The like remarks apply to the inversion $cdf$, $ce^2$.

\begin{center}
\textit{The MacMahon Linkage.} Art.\ Nos.\ 48 to 52.
\end{center}

48.\ We require the two theorems:

The first is: if a seminvariant $S$ has $q$ for its highest letter, then $\partial_q S$ is also a seminvariant.

The second has presented itself for unitariants (\textit{ante} No.\ 31); for seminvariants the form is less simple, viz.\ If in any seminvariant, attending only to the terms of the highest degree, we therein change $b, c, d, e, \ldots$ into $b, 2c, 6d, 24e, \ldots$ and then diminish the letters (that is, replace each letter by the next preceding letter) and in the result so obtained change $b, c, d, e, \ldots$ into $b,\ \tfrac{c}{2},\ \tfrac{d}{6},\ \tfrac{e}{24},\ \ldots$ we obtain a seminvariant. For instance $g - 6bf + 15ce - 10d^2$, in the terms of degree 2, making the numerical change we have $-720 bf + 720 ce - 360 d^2$, and then diminishing the letters and making the numerical change, we obtain $-720\,\dfrac{ae}{24} + 720\,\dfrac{bd}{6} - 360\,\dfrac{c^2}{4}$, that is, a seminvariant $-30(e - 4bd + 3c^2)$.

For the proof, observe that the equation $\Delta S = 0$, attending therein only to the terms of the highest degree, gives $(2b\partial_c + 3c\partial_d + \ldots) S' = 0$, if $S'$ denote the terms of the highest degree: making the numerical change, this is $(b\partial_c + c\partial_d + \ldots) S''$, if $S''$ is what $S'$ becomes thereby; diminishing the letters, this is $(\partial_b + b\partial_c + \ldots) S''' = 0$, if $S'''$
% ----- printed page 299 (PDF 321) -----
is the diminished value of $S''$, and finally making the numerical change, if $T$ be what $S'''$ becomes on writing therein $b, \tfrac{c}{2}, \tfrac{d}{6}, \ldots$ for $b, c, d, \ldots$, this gives
\[
(\partial_b + 2b\partial_c + \ldots) T = 0,
\]
viz.\ $T$ is a seminvariant.

49.\ Assume that, for the weights up to a certain weight $w$, the forms of the sharp seminvariants are known: and for the weight $w$ consider a seminvariant $I(\text{ao}) F$: here if $I$ be given, the first theorem establishes a limit $F'$ such that $F$ is in AO not higher than $F'$. For instance, when $w = 12$, if $I = dj$, the coefficient of $j$ as being a seminvariant can only be $d\,\mathrm{ao}\,b^3$, and thus the seminvariant contains a term $b^3 j$, or the final term $F$ must be in AO not higher than $b^3 j$; the degree is thus $= 4$ at least.

Similarly, if $F$ be given, then the second theorem determines a limit $I'$ such that $I$ is in CO not lower than $I'$. Thus when $w = 12$, as before, if $F = b^4 c d^2$, then diminishing the letters we have $bc^2$, a term belonging to $f\,\mathrm{ao}\,bc^2$; the diminished form has thus terms $a^4(a^2 f, bc^2)$, so that augmenting these the seminvariant has terms $bcg$ and $cd^2$, and thus the initial term $I$ is in CO not lower than $bcg$.

50.\ A limit for $I$ or $F$, when the other is given, can also in some cases be found as follows: Considering a seminvariant of the weight $w$ as before, and denoting its extent and degree by $\sigma$ and $\delta$ respectively, then we have $\sigma\delta - 2w\sigma = 0$ or positive; that is, $\sigma\delta = 2w$ at least; here given $I$, we have $\sigma$, and then $\delta = \dfrac{2w}{\sigma}$ at least: and given $F$ we have $\delta$, and then $\sigma = \dfrac{2w}{\delta}$ at least.

51.\ We may now explain the MacMahon linkage; for a given weight, we write down in two columns the initials or non-unitaries in CO, and the finals or power-enders in AO: by what precedes, it appears that we cannot combine the terms of the one column each with the term opposite to it in the other column; what we do is: beginning with the top of the column of initials, we combine successively each term with the highest admissible term in the column of finals: or beginning with the bottom of the column of finals, we combine successively each term with the lowest admissible term in the column of initials.

% ----- printed page 300 (PDF 322) -----
52.\ For the weight 12, the linkage is
\begin{center}
\textit{read downwards.}\\[2pt]
\begin{tabular}{r@{\;}c@{\;}l}
shown by not in AO higher than & & \\
$(c\,\mathrm{ao}\,b^2)k$ & $\longrightarrow$ & $bl$,\quad $k\,\mathrm{ao}\,l^2$ \\
$(d\,\mathrm{ao}\,b^3)j$ & $\longrightarrow$ & $b^2 k$,\quad $j\,\mathrm{ao}\,be^2$ \\
$(e\,\mathrm{ao}\,c^2)i$ & $\longrightarrow$ & $bdi$,\quad $ch\,\mathrm{ao}\,d^2$ \\
$(c^2\,\mathrm{ao}\,b^4)i$ & & $b^3 j$,\quad $i\,\mathrm{ao}\,e^2$ \\
$(f\,\mathrm{ao}\,bc^2)h$ & $\longrightarrow$ & $b^2 dh$,\quad $eg\,\mathrm{ao}\,cd^2$ \\
$(cd\,\mathrm{ao}\,b^5)h$ & & $b^4 i$,\quad $h\,\mathrm{ao}\,bd^2$ \\
$(g\,\mathrm{ao}\,d^2)g$ & $\longrightarrow$ & $b^4 c^2 e$,\quad $b^3 dg$,\quad $cf\,\mathrm{ao}\,bc^3$ \\
$(ce\,\mathrm{ao}\,c^3)g$ & & $b^{10}c$,\quad $b^3 ef$,\quad $de\,\mathrm{ao}\,b^3 c^2$ \\
$(d^2\,\mathrm{ao}\,b^2 c^4)g$ & & $b^5 h$,\quad $g\,\mathrm{ao}\,d^2$ \\
$(c^3\,\mathrm{ao}\,b^6)g$ & $\longrightarrow$ & $f^2$,\quad $cf\,\mathrm{ao}\,e^2$ \\
$(f^2\,\mathrm{ao}\,b^4 c^2)f$ & & $b^4 g$,\quad $f\,\mathrm{ao}\,b^2 d^2$ \\
$(de\,\mathrm{ao}\,b^3 c^2)f$ & & $b^2 g$,\quad $e\,\mathrm{ao}\,b c^2$ \\
$(c^2 d\,\mathrm{ao}\,b^7)f$ & $\longrightarrow$ & $b^8 g$,\quad $f\,\mathrm{ao}\,b^2 c^2$ \\
$(e^2\,\mathrm{ao}\,c^4)e$ & & $b^5 d c^7$,\quad $cd\,\mathrm{ao}\,b^5 c f^2$ \\
$(c^2 e\,\mathrm{ao}\,b^4 c^3)e$ & & $b^7 e$,\quad $d\,\mathrm{ao}\,b^3$ \\
$(c d^2\,\mathrm{ao}\,b^4 c^2)e$ & & $b^9 d$,\quad $c\,\mathrm{ao}\,b^2$ \\
$(c^4\,\mathrm{ao}\,b^8)e$ & & lower than shown by not in CO \\
$(d^3\,\mathrm{ao}\,b^4 c^2)d$ & & \\
$(c^2 d\,\mathrm{ao}\,b^9)d$ & & \\
\end{tabular}\\[2pt]
\textit{read upwards.}
\end{center}

Thus, beginning at the top of the column of initials, $m$ is to be linked with $g^2$, that is, we have $(m\,\mathrm{ao}\,g^2)$; $ck$ with $cf^2$, that is, we have $(ck\,\mathrm{ao}\,cf^2)$; $dj$ cannot be linked with $e^3$, for the final must be in AO not higher than $b^3 j$, but it is linked with the highest term $b^2 f^2$ for which this condition is satisfied, that is, we have $(dj\,\mathrm{ao}\,b^2 f^2)$; $ei$ is then linked with the highest admissible term $e^3$, that is, we have $(ei\,\mathrm{ao}\,e^3)$; and so on.

% ----- printed page 301 (PDF 323) -----
Or beginning at the bottom of the column of finals, $b^{12}$ is linked with $c^6$, that is, we have $(c^6\,\mathrm{ao}\,b^{12})$, $b^8 c^2$ with $c^3 d^2$, that is, we have $(c^3 d^2\,\mathrm{ao}\,b^8 c^2)$; $b^6 c^3$ cannot be linked with $d^4$, for the initial must be in CO not lower than $b^8 e$, but it is linked with the lowest term $c^4 e$ for which this condition is satisfied, that is, we have $(c^4 e\,\mathrm{ao}\,b^6 c^3)$; and so on.

\begin{center}
\textit{The Umbral Notation. Stroh's Theory.} Art.\ Nos.\ 53 to 56.
\end{center}

53.\ Employing the umbrae $\alpha$, $\beta$, $\gamma$, $\delta$, \ldots, which are such that
\[
\alpha = \beta = \gamma = \ldots = b;\quad \alpha^2 = \beta^2 = \gamma^2 = \ldots = c;\quad \alpha^3 = \beta^3 = \gamma^3 = \ldots = d;
\]
and so on, then for instance
\[
(\alpha - \beta)^2 = \alpha^2 - 2\alpha\beta + \beta^2 = c - 2b^2 + c = 2(c - b^2),
\]
a seminvariant;
\[
(\alpha - \beta)^2(\alpha - \gamma) = \alpha^3 - 2\alpha^2\beta + \alpha\beta^2 - \alpha^2\gamma + 2\alpha\beta\gamma - \beta^2\gamma,
\]
\[
= d - 2bc + b\cdot c - bc + 2b^3 - bc = d - 3bc + 2b^3,
\]
a seminvariant: and so in general any rational and integral function of the differences of the umbrae developed and interpreted is a seminvariant. For the seminvariants of a given weight, e.g.\ $w = 6$, Dr Stroh\,$^*$ considers the function
\[
\Omega^6 = (\alpha x + \beta y + \gamma z + \delta w + \epsilon t + \zeta u)^6,
\]
where $x, y, z, w, t, u$ are numbers the sum of which is $= 0$, or we may if we please have more than 6 such numbers: the expression is obviously a function of the differences of the umbrae and it is thus a seminvariant. To develop its value, observe that after expansion of the sixth power we have sets of similar terms, for instance $\alpha^6 + \beta^6 + \gamma^6 + \ldots$ which putting therein $\alpha^6 = \beta^6 = \gamma^6 = \ldots = g$ become $= g\cdot S x^6$, and generally each set becomes equal to a literal term multiplied by a symmetric function of the $x, y, z, w, \ldots$; introducing capital letters to denote the elementary symmetric functions of these quantities, and recollecting that their sum is assumed to be $= 0$, say we have
\[
1 + C u^2 + D u^3 + E u^4 + \ldots = \overline{1 - xu}\cdot\overline{1 - yu}\cdot\overline{1 - zu}\cdots,
\]
(that is, $0 = S x$, $+C = S xy$, $-D = S xyz$, \&c.) then by aid of the Table VI ($b$) writing therein $0, C, D, E, F, G$ for $b, c, d, e, f, g$, we find
\[
\Omega^6 = (\alpha x + \beta y + \gamma z + \ldots)^6 = g\cdot S x^6 + 6 S\alpha^5\beta\cdot S x^5 y + \&c.,
\]

\noindent\footnotesize $^*$ See the paper ``Ueber die Symbolische Darstellung der Grundsyzyganten einer bin\"aren Form sechster Ordnung und eine Erweiterung der Symbolik von Clebsch,'' \textit{Math.\ Ann.}\ t.\ \textsc{xxxvi}.\ (1890), pp.\ 262--303, in particular \S 10, \textit{Das Formensystem einer Form unbegrenzt hoher Ordnung}.\normalsize

% ----- printed page 302 (PDF 324) -----
as shown in the following table:

\begin{center}
\begin{tabular}{l|rrrr}
 & $C^3$ & $D^2$ & $CE$ & $G$ \\
\hline
$1\cdot g$ & $-2$ & $-3$ & $+6$ & $-6$ \\
$6\cdot bf$ & $+2$ & $-3$ & $-6$ & $+6$ \\
$15\cdot ce$ & $-2$ & $-3$ & $+2$ & $+6$ \\
$20\cdot d^2$ & $+1$ & $+3$ & $-3$ & $+3$ \\
$30\cdot b^2 e$ & & $-3$ & $-2$ & $-6$ \\
$60\cdot bcd$ & & $-3$ & $+4$ & $-12$ \\
$90\cdot c^3$ & & $+1$ & $-2$ & $-2$ \\
$120\cdot b^3 d$ & & & $-2$ & $+6$ \\
$180\cdot b^2 c^2$ & & & $+1$ & $+9$ \\
$360\cdot b^4 c$ & & & & $-6$ \\
$720\cdot b^6$ & & & & $+1$ \\
\hline
& $[d^2]$ & $[c^3]$ & $[b^2 c^2]$ & $[b^6]$ \\
\end{tabular}
\end{center}

the numbers whereof are, it will be observed, identical with those of the foregoing table No.\ 33, relating to the MacMahon equation.

This is to be read
\[
\Omega^6 = C^3 [d^2] + D^2 [c^3] + CE [b^2 c^2] + G [b^6],
\]
viz.\ $\Omega^6$ is a linear function of $C^3$, $D^2$, $CE$ and $G$, the coefficients of these, being given functions of $(b, c, d, e, f, g)$, which given functions are the specific blunt seminvariants which have been already called $[d^2]$, $[c^3]$, $[b^2 c^2]$ and $[b^6]$. And so in general, the developed value of $\Omega^w$ affords a complete definition of these specific blunt seminvariants of the weight $w$. Observe that $\alpha, \beta, \gamma, \delta, \ldots$ are umbrae in nowise connected with the roots $\alpha, \beta, \gamma, \delta, \ldots$ before made use of, and that $B, C, D, \ldots$ are actual quantities in nowise connected with the symbolic capitals $B, C, D, \ldots$ before made use of.

54.\ The capital and small letter symbols are conjugate to each other. It will be convenient to give here, in reference to subsequent investigations, a table of these conjugate forms up to the degree 6 and weight 15.

% ----- printed page 303 (PDF 325) -----
\textit{[Figure: a stepped, page-wide ``Table of Conjugates'' giving capital-letter and small-letter conjugate forms for every weight from 1 (left column) up to 15 (right edge) and for degrees 2 through 6 (rows arranged in a step pattern from upper right downwards to lower left). Each cell pairs a small-letter monomial (e.g.\ $bf$, $ce$, $d^2$, $b^2 e$, \ldots) with its capital-letter conjugate (e.g.\ $BF$, $CE$, $D^2$, $B^2 E$, \ldots), in the staircase layout of the original. The detailed cell-by-cell typesetting is too damaged at scan resolution to reproduce literally; the principle of construction is that for every weight $w$ and every degree $\delta$ ($2 \leq \delta \leq 6$, $w \leq 15$), each non-unitary monomial in the symmetric functions of weight $w$ and degree $\delta$ appears together with its symbolic-capital conjugate, in agreement with the convention introduced in Art.\ 53 above.]}

\end{document}
